% Compile with: uplatex file.tex && dvipdfmx file.dvi
%           or: latexmk -uplatex file.tex
\documentclass[dvipdfmx]{article}
\usepackage{amsmath, amssymb}
\usepackage{tikz}

\title{scenario\_kspace\_kinetic\_term}
\date{\today}

\begin{document}
\maketitle

Check that the most ordinary condensed-matter kinetic term $\xi_k c^\dagger_k c_k$ can be written. $\xi_k$ is one quantity ("the kinetic energy at $k$"), so its subscript is declared as part of the name, not as a component index.

c as the annihilation operator for electron

\[
  c
\]

Define $\xi_{k}$ as a scalar

\[
  \xi_{k}
\]

Define Delta as a scalar

\[
  \Delta
\]

Write the kinetic term. An indexed scalar multiplying QA operators reaches the CAS as written.

$\xi_{k}$ c'(k) c(k)

\begin{equation}
  \xi_{k} c^{\dagger}(k) c(k)
  \label{eq:E1}
\end{equation}

Normal order it. The expression is already ordered, so nothing changes — the point is that the display keeps the researcher's own subscript $\xi_k$ and no CAS syntax leaks out.

$normal_order$ E1

\begin{equation}
  \xi_{k} {c}_{k}^{\dagger} {c}_{k}
  \label{eq:E2}
\end{equation}

Now a Hamiltonian with the product deliberately reversed, plus a pair term to commute with later.

$H = \xi_{k} c(k) c'(k) + \Delta c'(k) c'(-k)$

\begin{equation}
  H = \xi_{k} c(k) c^{\dagger}(k) + \Delta c^{\dagger}(k) c^{\dagger}(-k)
  \label{eq:E3}
\end{equation}

Normal order it. The anticommutator $c c^\dagger = 1 - c^\dagger c$ produces the constant term $\xi_k$ — evidence that the indexed scalar passes through the anticommutator intact.

$normal_order$ E3

\begin{equation}
  H = \xi_{k} -  \xi_{k} c^{\dagger}_{k} c_{k} + \Delta c^{\dagger}_{k} c^{\dagger}_{-k}
  \label{eq:E4}
\end{equation}

Take the commutator $[H, c_k]$, the backbone of the equation of motion. By hand $[c^\dagger_k c_k, c_k]=-c_k$ and $[c^\dagger_k c^\dagger_{-k}, c_k]=-c^\dagger_{-k}$, so the answer should be $\xi_k c_k - \Delta c^\dagger_{-k}$.

commutator of E3 and c(k)

\begin{equation}
  - \Delta c^{\dagger}_{-k} + \xi_{k} c_{k}
  \label{eq:E5}
\end{equation}

\end{document}